\section{收敛速度}
\label{sec:graph-walk-convergence}

上一节末尾看到，在二部图上单点出发的随机游走会永久奇偶振荡。对于一般的连通无向图，从任意初态$\bm\nu^{(0)}$出发的多步分布$\bm\nu^{(t)}=\bm\nu^{(0)}\bm P^t$究竟何时收敛到平稳分布$\bm\pi$？如果收敛，需要走多少步才能与平稳分布足够接近？本节先将转移矩阵$\bm P$对称化以对接归一化拉普拉斯谱，再导出收敛的充要条件与误差衰减界，最后对比几类典型图的收敛速度。

\subsection{转移矩阵的对称化}
\label{sec:graph-walk-symmetrization}

在度数不均匀的无向图上，虽然边权对称（$w_{i,j}=w_{j,i}$），转移矩阵$\bm P=\bm D^{-1}\bm W$本身通常并\emph{不对称}（例如在四点路径上，$P_{1,2}=1$而$P_{2,1}=1/2$）。为了使用实对称矩阵的正交谱分解，我们对$\bm P$作度平方根相似变换：将第$i$行乘以$\sqrt{d_i}$、第$j$列除以$\sqrt{d_j}$，其$(i,j)$元变为对称的$\frac{w_{i,j}}{\sqrt{d_i d_j}}$。写成矩阵形式，即左乘$\bm D^{1/2}$、右乘$\bm D^{-1/2}$：
\begin{equation}
\bm{S}=\bm{D}^{1/2}\bm{P}\bm{D}^{-1/2}
=\bm{D}^{-1/2}\bm{W}\bm{D}^{-1/2}=\bm{I}-\mathcal{L},
\label{eq:graph-walk-symmetric-transition}
\end{equation}
其中$\mathcal L=\bm D^{-1/2}\bm L\bm D^{-1/2}$正是第~\ref{sec:graph-normalized-laplacian}节的对称归一化拉普拉斯矩阵。

式~\eqref{eq:graph-walk-symmetric-transition}建立起了随机游走与归一化拉普拉斯谱之间的直接桥梁：
\begin{enumerate}
\item \emph{特征值一一对应}：因为$\bm P=\bm D^{-1/2}\bm S\bm D^{1/2}$与实对称矩阵$\bm S$相似，两者的特征值完全相同。设$\mathcal L$的特征值按升序排列为$0=\mu_1<\mu_2\leq\cdots\leq\mu_n\leq 2$，则$\bm P$与$\bm S$的特征值按降序排列为
\[
  1=\theta_1>\theta_2\geq\cdots\geq\theta_n\geq -1,
  \qquad \theta_i = 1 - \mu_i.
\]
\item \emph{多步演化的对称坐标表示}：由$\bm P\bm D^{-1/2}=\bm D^{-1/2}\bm S$，在多步演化式$\bm\nu^{(t+1)}=\bm\nu^{(t)}\bm P$两侧右乘$\bm D^{-1/2}$，可得
\[
\bm{\nu}^{(t+1)}\bm{D}^{-1/2}
=(\bm{\nu}^{(t)}\bm{D}^{-1/2})\bm{S}
=(\bm{\nu}^{(0)}\bm{D}^{-1/2})\bm{S}^{t+1}.
\]
也就是说，只要把概率分布行向量$\bm\nu^{(t)}$缩放为$\bm\nu^{(t)}\bm D^{-1/2}$，每走一步就等价于右乘一次实对称矩阵$\bm S$。
\end{enumerate}

例如在单位权四点路径上，$\bm P$与$\bm S$的四个特征值为$\{1,\frac12,-\frac12,-1\}$；在单位权四元环$C_4$上，特征值为$\{1,0,0,-1\}$。两者的最小特征值都达到了$-1$（对应$\mu_4=2$）。

\subsection{收敛条件}
\label{sec:graph-walk-convergence-condition}

利用对称矩阵$\bm S$的标准正交特征基，可以一步写出第$t$步分布$\bm\nu^{(t)}$与平稳分布$\bm\pi$之间的偏差演化公式。

设$\bm u_1,\bm u_2,\ldots,\bm u_n$为$\bm S$对应于$1=\theta_1>\theta_2\geq\cdots\geq\theta_n\geq -1$的标准正交特征向量（列向量），其中第一特征向量为$\bm u_1=\bm D^{1/2}\bm 1/\sqrt{\operatorname{vol}(V)}$。记第$t$步的分布偏差行向量为$\bm z^{(t)}=\bm\nu^{(t)}-\bm\pi$。由于$\bm\nu^{(t)}\bm 1=1$且$\bm\pi\bm 1=1$，偏差的分量和恒为零：$\bm z^{(t)}\bm 1=0$。因此在对称坐标下，缩放偏差$\bm z^{(t)}\bm D^{-1/2}$与第一特征向量$\bm u_1$恒正交：
\[
(\bm z^{(t)}\bm D^{-1/2})\bm u_1
=\frac{1}{\sqrt{\operatorname{vol}(V)}}(\bm z^{(t)}\bm D^{-1/2})(\bm D^{1/2}\bm 1)
=\frac{\bm z^{(t)}\bm 1}{\sqrt{\operatorname{vol}(V)}}=0.
\]
这意味着偏差在$\bm u_1$方向上的分量恒为$0$。将初始缩放偏差在其余$n-1$个非平稳特征方向$\bm u_2,\ldots,\bm u_n$上展开为$\bm z^{(0)}\bm D^{-1/2}=\sum_{i=2}^n a_i\bm u_i^\top$（其中$a_i=\bm z^{(0)}\bm D^{-1/2}\bm u_i$），每走一步右乘$\bm S$就将第$i$项乘上$\theta_i$，从而第$t$步的缩放偏差精确等于
\begin{equation}
\bm{z}^{(t)}\bm{D}^{-1/2}
=\sum_{i=2}^n a_i\theta_i^t\bm{u}_i^\top.
\label{eq:graph-walk-spectral-evolution}
\end{equation}

\begin{theorem}[无向连通图随机游走的收敛判别]
\label{thm:graph-walk-convergence}
设有限无向图的正权支撑连通且至少有两个顶点，平稳分布为$\bm\pi$。则从任意初始分布$\bm\nu^{(0)}$出发的随机游走分布$\bm\nu^{(t)}=\bm\nu^{(0)}\bm P^t$当$t\to\infty$时收敛于$\bm\pi$，当且仅当正权支撑不是二部图（等价于最小特征值$\theta_n>-1$，或$\mu_n<2$）。
\end{theorem}

\begin{proof}
由式~\eqref{eq:graph-walk-spectral-evolution}及$\bm u_2,\ldots,\bm u_n$的标准正交性，有$\lVert\bm z^{(t)}\bm D^{-1/2}\rVert_2^2=\sum_{i=2}^n a_i^2\theta_i^{2t}$。该平方和在$t\to\infty$时对任意初始系数$(a_2,\ldots,a_n)$趋于$0$，当且仅当对所有$i=2,\ldots,n$都有$|\theta_i|<1$。由于连通图上$1>\theta_2\geq\cdots\geq\theta_n\geq -1$，故$\max_{i\geq 2}|\theta_i|<1$当且仅当$\theta_n>-1$（即$\mu_n=1-\theta_n<2$）。由第~\ref{sec:graph-normalized-laplacian}节定理~\ref{prop:graph-normalized-spectrum-bound}，这等价于图不是二部图。必要性并不要求任意系数组合都能表示概率分布：若存在$\theta_i=-1$，可取$\bm\nu^{(0)}=\bm\pi+a\bm u_i^\top\bm D^{1/2}$，其中非零$a$足够小。其分量和为$1$且逐项非负，却保留不衰减的交替模态。
\end{proof}

图~\ref{fig:graph-walk-parity}直观对比了非二部图与二部图的演化差异：在三角形$C_3$（非二部图，$\theta_2=\theta_3=-1/2$）上，偏差按$(-1/2)^t$指数衰减，分布迅速趋向均匀平稳分布；而在四元环$C_4$（二部图，$\theta_4=-1$）上，$\theta_4=-1$对应的正负交替模态$(-1)^t\bm u_4^\top$永不衰减，概率质量在奇偶两侧顶点集之间无休止地来回跳动。

\begin{figure}[!htb]
\centering
% Exact transient probabilities on the unweighted triangle and four-cycle.
\begin{tikzpicture}[x=1mm,y=1mm,font=\figurefont\fontsize{8}{9}\selectfont]
  \foreach \column/\step in {0/0,1/1,2/2,3/3} {
    \node[text=FigureInk] at ({20*\column+8},35) {$t=\step$};
  }
  \node[anchor=east,text=FigureInk] at (-2,25) {$C_3$};
  \node[anchor=east,text=FigureInk] at (-2,7) {$C_4$};
  \foreach \column/\values in {0/{12,0,0},1/{0,6,6},2/{6,3,3},3/{3,4.5,4.5}} {
    \begin{scope}[shift={({20*\column},19)}]
      \foreach \value [count=\vertex from 1] in \values {
        \fill[FigureBlue] ({4*\vertex-2},0) rectangle ({4*\vertex+1},\value);
        \node[anchor=north,text=FigureInk] at ({4*\vertex-.5},0) {$\vertex$};
      }
      \draw[FigureStroke,line width=.55pt] (0,0)--(17.5,0);
      \draw[FigureStroke,dashed,line width=.55pt] (0,4)--(17.5,4);
    \end{scope}
  }
  \foreach \column/\values in {0/{12,0,0,0},1/{0,6,0,6},2/{6,0,6,0},3/{0,6,0,6}} {
    \begin{scope}[shift={({20*\column},1)}]
      \foreach \value [count=\vertex from 1] in \values {
        \fill[FigureBlue] ({4*\vertex-2},0) rectangle ({4*\vertex+1},\value);
        \node[anchor=north,text=FigureInk] at ({4*\vertex-.5},0) {$\vertex$};
      }
      \draw[FigureStroke,line width=.55pt] (0,0)--(17.5,0);
      \draw[FigureStroke,dashed,line width=.55pt] (0,3)--(17.5,3);
    \end{scope}
  }
\end{tikzpicture}

\caption{从顶点$1$出发的前四个时刻位置分布对比。每个面板自左至右以柱高表示各顶点的概率，水平虚线为均匀平稳概率。三角形$C_3$（上排）按$(-1/2)^t$迅速衰减收敛；四元环$C_4$（下排）因含$\theta_4=-1$模态，质量始终在二部划分的两侧之间交替跳动。}
\label{fig:graph-walk-parity}
\end{figure}

\subsection{收敛速度}
\label{sec:graph-walk-lazy}

当随机游走满足收敛条件时，非平稳特征值的最大模长
\[
  \rho = \max_{2\leq i\leq n}|\theta_i| = \max\{\theta_2, -\theta_n\} < 1
\]
直接决定了偏差$\bm z^{(t)}=\bm\nu^{(t)}-\bm\pi$向零衰减的几何速率，而$1-\rho>0$则称为转移矩阵的\term{绝对谱隙}（absolute spectral gap）。

\paragraph{总变差距离与谱间隙衰减界}

为了定量衡量第$t$步分布$\bm\nu^{(t)}$与平稳分布$\bm\pi$的接近程度，概率论中标准采用\term{总变差}（total variation）距离：
\[
  \lVert\bm\nu-\bm\pi\rVert_{\mathrm{TV}}
  =\frac12\lVert\bm\nu-\bm\pi\rVert_1
  =\frac12\sum_{i=1}^n|\nu_i-\pi_i|
  =\max_{A\subseteq V}\bigl|\bm\nu(A)-\bm\pi(A)\bigr|.
\]
另一方面，在对称坐标式~\eqref{eq:graph-walk-spectral-evolution}中，最契合正交展开的欧氏范数是以平稳概率$\pi_i=d_i/\operatorname{vol}(V)$倒数加权的$\pi^{-1}$-范数：
\[
\lVert\bm{z}\rVert_{\pi^{-1}}^2
=\sum_{i=1}^n\frac{z_i^2}{\pi_i}
=\operatorname{vol}(V)\,\lVert\bm{z}\bm{D}^{-1/2}\rVert_2^2.
\]
下面的定理通过Cauchy--Schwarz不等式将两者衔接，给出总变差距离随步数$t$的指数衰减上界。

\begin{theorem}[谱间隙与总变差收敛速度]
\label{prop:graph-walk-tv-bound}
设连通正权无向图随机游走的非平稳特征值最大模长为$\rho=\max_{i\geq 2}|\theta_i|$。则对任意初始分布$\bm\nu^{(0)}$及任意步数$t\geq 0$，有
\begin{equation}
\lVert\bm{\nu}^{(t)}-\bm{\pi}\rVert_{\pi^{-1}}
\le\rho^t\lVert\bm{\nu}^{(0)}-\bm{\pi}\rVert_{\pi^{-1}},
\qquad
\lVert\bm{\nu}^{(t)}-\bm{\pi}\rVert_{\mathrm{TV}}
\le\frac12\rho^t\lVert\bm{\nu}^{(0)}-\bm{\pi}\rVert_{\pi^{-1}}.
\label{eq:graph-walk-weighted-contraction}
\end{equation}
特别地，若从单个顶点$u\in V$确定性出发（$\bm\nu^{(0)}=\bm e_u^\top$），则
\begin{equation}
\lVert\bm{\nu}^{(t)}-\bm{\pi}\rVert_{\mathrm{TV}}
\le\frac12\sqrt{\frac{1}{\pi_u}-1}\;\rho^t.
\label{eq:graph-walk-tv-bound}
\end{equation}
\end{theorem}

\begin{proof}
记$\bm z^{(t)}=\bm\nu^{(t)}-\bm\pi$。由正交展开式~\eqref{eq:graph-walk-spectral-evolution}，有
\[
\lVert\bm{z}^{(t)}\bm{D}^{-1/2}\rVert_2^2
=\sum_{i=2}^n a_i^2|\theta_i|^{2t}
\le \rho^{2t}\sum_{i=2}^n a_i^2
=\rho^{2t}\lVert\bm{z}^{(0)}\bm{D}^{-1/2}\rVert_2^2.
\]
两边同乘$\operatorname{vol}(V)$并开方即得式~\eqref{eq:graph-walk-weighted-contraction}第一式。接着将$|z_i^{(t)}|$分解为$\sqrt{\pi_i}\cdot\frac{|z_i^{(t)}|}{\sqrt{\pi_i}}$，应用Cauchy--Schwarz不等式并利用$\sum_{i=1}^n\pi_i=1$，得
\[
2\lVert\bm z^{(t)}\rVert_{\mathrm{TV}}
=\sum_{i=1}^n|z_i^{(t)}|
\le\left(\sum_{i=1}^n\pi_i\right)^{1/2}
   \left(\sum_{i=1}^n \frac{(z_i^{(t)})^2}{\pi_i}\right)^{1/2}
=\lVert\bm{z}^{(t)}\rVert_{\pi^{-1}}
\le\rho^t\lVert\bm{z}^{(0)}\rVert_{\pi^{-1}},
\]
同除以$2$即得式~\eqref{eq:graph-walk-weighted-contraction}第二式。当$\bm\nu^{(0)}=\bm e_u^\top$时，$z_u^{(0)}=1-\pi_u$，对$j\neq u$有$z_j^{(0)}=-\pi_j$，故$\lVert\bm z^{(0)}\rVert_{\pi^{-1}}^2=\frac{(1-\pi_u)^2}{\pi_u}+(1-\pi_u)=\frac{1}{\pi_u}-1$，代入即得式~\eqref{eq:graph-walk-tv-bound}。
\end{proof}

由式~\eqref{eq:graph-walk-tv-bound}可知，要使总变差距离降低到给定精度$0<\varepsilon<1/2$以内，利用$\rho^t=(1-(1-\rho))^t\leq\mathrm e^{-(1-\rho)t}$，定义\term{混合时间}（mixing time）为使所有单点初态的总变差误差均不超过$\varepsilon$的最小非负整数时刻。记$\pi_{\min}=\min_i\pi_i$，在$\rho<1$时，上述界给出
\begin{equation}
  t_{\mathrm{mix}}(\varepsilon) \leq \left\lceil\frac{1}{1-\rho}\log\frac{1}{2\varepsilon\sqrt{\pi_{\min}}}\right\rceil.
  \label{eq:graph-walk-mixing-time-bound}
\end{equation}
这一充分步数上界包含$(1-\rho)^{-1}$，同时依赖目标误差与$\pi_{\min}$；它不是实际混合时间与谱隙的精确反比关系。

\paragraph{惰性游走：消除负特征值振荡}

在二部图或接近二部图的网络中，最小特征值$\theta_n\approx -1$会导致$\rho=\max\{\theta_2,-\theta_n\}$由高频奇偶振荡项$-\theta_n$主导。为使转移矩阵的全部特征值非负，算法应用中常采用\term{惰性游走}（lazy random walk）：每一步先抛一枚公平硬币，以$1/2$的概率停留在当前顶点不动，以$1/2$的概率按原转移矩阵$\bm P$走一步。惰性游走的转移矩阵为
\begin{equation}
\bm{Q}=\frac{\bm{I}+\bm{P}}{2}=\bm I - \frac12\bm L_{\mathrm{rw}}.
\label{eq:graph-walk-lazy-transition}
\end{equation}
由于原地停留和按$\bm P$转移都保持$\bm\pi$不变，惰性游走具有完全相同的平稳分布$\bm\pi$；但它的特征值被平移缩放为$\theta_i(\bm Q)=\frac{1+\theta_i}{2}=1-\frac{\mu_i}{2}\in[0,1]$，全部变为非负实数。此时原先$\theta_n=-1$的奇偶振荡模态直接变为$0$，非平稳特征值的最大模长完全由归一化拉普拉斯的第二小特征值$\mu_2$决定：
\[
  \rho(\bm Q) = 1 - \frac{\mu_2}{2},
  \qquad
  \text{惰性谱隙 } \gamma = 1 - \rho(\bm Q) = \frac{\mu_2}{2}.
\]

\subsection{典型图的收敛速度}
\label{sec:graph-walk-bottleneck}

将上一小节的惰性谱隙$\gamma=\mu_2/2$与第~\ref{sec:graph-cheeger}节的Cheeger不等式$\frac{\mu_2}{2}\leq\phi_*\leq\sqrt{2\mu_2}$结合，立即可得谱隙$\gamma$与全图最小导通率$\phi_*$的双向控制关系：
\begin{equation}
\frac{\phi_*^2}{4}\le\gamma\le\phi_*.
\label{eq:graph-walk-lazy-cheeger}
\end{equation}
式~\eqref{eq:graph-walk-lazy-cheeger}从几何上解释了什么图收敛快、什么图收敛慢：当$\operatorname{vol}(S)\leq\operatorname{vol}(V)/2$时，普通游走的导通率$\phi(S)=\frac{c(S,\bar S)}{\operatorname{vol}(S)}$正是平稳状态下棋子位于集合$S$内时“下一步跨过边界离开$S$的条件概率”；如果图中存在导通率$\phi_*$极小的稀疏割瓶颈，棋子一旦落入瓶颈一侧就很难穿过狭窄的咽喉逃到另一侧，谱隙$\gamma\leq\phi_*$必然极小，混合时间随之急剧拉长。

下面固定总变差误差$0<\varepsilon<1/2$，比较四类图族的谱隙与混合时间。除明确说明的普通游走外，均使用惰性游走；典型量级中的常数可以依赖$\varepsilon$。

\paragraph{1. 完全图$K_n$（固定误差下常数步混合，$t_{\mathrm{mix}}=\Theta(1)$）}

在单位权完全图$K_n$上（$n\geq3$），归一化拉普拉斯的所有非零特征值均为$\mu_2=\cdots=\mu_n=\frac{n}{n-1}>1$。普通游走的非平稳特征值全为$\theta_i=-\frac{1}{n-1}$，谱半径$\rho=\frac{1}{n-1}$随$n$增大反而趋于$0$；即便采用惰性游走，谱隙也高达$\gamma=\frac{n}{2(n-1)}>\frac12$。从任意顶点出发，仅需$O(1)$步即可在整个完全图上接近均匀分布。

\paragraph{2. 超立方体$Q_d$（对数级快速混合，$t_{\mathrm{mix}}=\Theta(d\log d)=\Theta(\log n\log\log n)$）}

在顶点数$n=2^d$的$d$维超立方体$Q_d$上，最小导通率为$\phi_*=1/d$，归一化第二小特征值为$\mu_2=2/d$（第~\ref{sec:graph-cheeger-tightness}小节）。惰性游走的谱隙为$\gamma=\frac{1}{d}=\frac{1}{\log_2 n}$。代入式~\eqref{eq:graph-walk-mixing-time-bound}（其中$\pi_u=2^{-d}$给出$\log(1/\sqrt{\pi_u})=\Theta(d)$），可知混合时间为$O(d^2)$（进一步利用坐标更新过程的耦合与谱展开可证明$t_{\mathrm{mix}}=\Theta(d\log d)$），虽然顶点数呈指数增长，随机游走却能在多项式对数步内迅速混合。

\paragraph{3. 路径图$P_n$与环图$C_n$（二次方扩散混合，$t_{\mathrm{mix}}=\Theta(n^2)$）}

在$n$点路径$P_n$或环$C_n$上，最小导通率为$\phi_*=\Theta(1/n)$，而归一化第二小特征值为$\mu_2=1-\cos(\Theta(1/n))=\Theta(n^{-2})$，惰性谱隙仅为$\gamma=\Theta(n^{-2})$。远离端点时，游走的位移呈扩散尺度，经过$t$步的典型距离为$O(\sqrt t)$，这解释了跨越长度$n$为何需要二次方时间。端点反射和惰性停留需另行计入；严格的$\Theta(n^2)$混合界可由路径、环的转移谱展开得到，不能只由单步方差推出。

\paragraph{4. 杠铃图（细腰瓶颈严重拖慢混合，$t_{\mathrm{mix}}=\Omega(n^2)$）}

将两个各含$k=n/2$个顶点的完全图$K_k$用一条单位权桥边相连（图~\ref{fig:graph-connectivity-archetypes}b）。虽然团内的位置分布在常数步内就可接近均匀，但单侧团的体积高达$\operatorname{vol}(S)=k(k-1)+1=\Theta(n^2)$，而跨桥割权仅为$c(S,\bar S)=1$，导致最小导通率骤降至
\[
  \phi_* = \frac{1}{k(k-1)+1} = \Theta(n^{-2}).
\]
由式~\eqref{eq:graph-walk-lazy-cheeger}的上界$\gamma\leq\phi_*=\Theta(n^{-2})$可知，这条细腰单桥构成了严重的扩散瓶颈：从团内平稳条件分布出发跨越桥边的等待具有$\Theta(n^2)$量级；谱隙的倒数也为混合时间提供相应下界，从而固定误差下全局混合时间为$\Omega(n^2)$。
